% READ-RECOMPUTE-Q4-8
\section*{READ-RECOMPUTE-Q4-8: STORE/RECOMP tradeoffs for the $Q_4$ edge-cover predicate}

\textbf{STATUS: PROVED/EXACT constructions plus COMPUTATIONAL finite node/path audits. Global minima remain OPEN. P versus NP remains OPEN.}

\paragraph{Corpus reconciliation.}
This run reviewed root \texttt{main.tex}, the repository research protocol, ADAPTIVE-FORMULA-FBDD5, READ-TWICE-Q4-6, and READ-REPETITION-Q4-7. READ-TWICE-Q4-6 already contains the exact 64-node syntactic read-twice construction. READ-REPETITION-Q4-7 contains a much looser greedy upper-bound curve and its claimed source/report provenance is incomplete in the current repository. The present run supersedes that curve as an upper bound, while preserving the earlier 64-node theorem.

\paragraph{Function.}
For full genuine 3-SAT on four Boolean variables, the 32 clause-presence variables identify with the 32 edges of $Q_4$. UNSAT is exactly the edge-cover predicate
\[
f(x)=\bigwedge_{v\in V(Q_4)}\bigvee_{e\ni v}x_e.
\]

\paragraph{A 64-node semantic $(1,+14)$ construction.}
The vertex-star program scans the four incident edge bits of each cube vertex until it finds a one. It has exactly
\[16\cdot4=64\]
query nodes and every input variable occurs globally in exactly two query nodes, so the program is syntactic read-twice.

A seeded local search over the four-edge order at each vertex followed by an exact consistent-path dynamic analysis gives
\[\max\{\text{total queries}\}=43\]
over the complete 32-bit input domain, and
\[\max\{\text{variables actually reread on one consistent path}\}=14.\]
Hence the same explicit 64-node program is a semantic $(1,+14)$-BP.

On the 272 perfect matchings of $Q_4$, every path uses exactly 40 queries and the repeated-variable count has
\[\min=8,\qquad \mathrm{avg}=10.117647\ldots,\qquad \max=13.\]
The lower endpoint has a direct proof: on a perfect matching, every selected matching edge is the unique selected incident edge at each endpoint, so all eight selected variables must be queried at both endpoints. Therefore every vertex-by-vertex star-scan implementation repeats at least eight variables on every perfect matching.

An independent exhaustive run over all
\[\binom{32}{8}=10{,}518{,}300\]
weight-eight inputs gives zero mismatches and a worst promised query count of 42.

\paragraph{Improved STORE/RECOMP curve.}
Fix vertex order $0,\ldots,15$ and ascending local edge-id order. Let $K$ be a set of edge variables whose first value is forgotten and may be queried again at the other endpoint. Values of variables outside $K$ are retained until their second endpoint and never reread. Identical future states are hash-consed exactly and redundant query nodes are removed. Thus every constructed program is syntactic read-twice and syntactic $(1,+|K|)$.

A greedy valid upper-bound curve is
\[
\begin{array}{c|rrrrrrrrrrrr}
k&0&1&2&4&8&12&16&20&24&28&29&30\\
\hline
|P|&1592&1193&945&610&298&168&114&88&70&65&64&62.
\end{array}
\]
This improves the previous endpoint-memory curve and answers the prior open question whether a hybrid can beat 64 nodes:
\[\boxed{|P_{\mathrm{read\mbox{-}twice}}|\le62}.\]

Exhausting all choices of stored variables at the high-recompute end gives the exact optimum \emph{inside this fixed endpoint-scan architecture}:
\[
\begin{array}{c|ccccc}
|K|&28&29&30&31&32\\
\hline
\min |P|&64&63&62&63&64.
\end{array}
\]
For $|K|=30$, the two stored variables are edges $(0,1)$ and $(14,15)$. Forcing additional forgetting can destroy useful merging. Under a budget of at most $k$ reread-enabled variables, the 62-node $k=30$ program remains available for $k=31,32$.

\paragraph{Relation to the lower-bound literature.}
Jukna--Razborov prove strong $(1,+s)$ lower bounds for suitable linear-code characteristic functions, while Savicky--Zak and Woelfel prove strict $(1,+k)$ hierarchies and mixedness-based lower bounds. Woelfel's theorem has the built-in condition $k\ell\le n$. For this 32-input laboratory, a lower bound on semantic $(1,+14)$ via that theorem would use $k=15$, hence $\ell\le2$ and can yield at most the constant bound $2^{\ell/2}\le2$. Thus that general mixedness theorem is parametrically incapable of certifying a meaningful lower bound near 64 in this high-reread regime.

Razgon's recent bounded-repetition work gives a direct conceptual match: rereading provides splitting power for CNFs whose constraint blocks can be represented separately. Our finite experiment is a small explicit STORE/FORGET/RECOMPUTE instance of the same phenomenon.

\paragraph{Asymptotic hostile control.}
The $Q_4$ read-twice collapse does not generalize directly. For $n$-variable genuine 3-SAT, the assignment-cover CNF has $2^n$ assignment constraints of width $\binom n3$, and every clause-presence variable occurs
\[2^{n-3}\]
times. The direct sequential construction therefore has
\[2^n\binom n3\]
query nodes and is syntactic read-$2^{n-3}$. At $n=5$ this is already a 320-node read-4 program. Hence the finite $n=4$ result is a hostile control, not an asymptotic SAT algorithm.

\paragraph{Barrier status.}
Rereading/recomputation is now an explicit escape from read-once state lower bounds. The project must charge a combined resource---retained state, reread work, or a time-space-recomputation tradeoff---rather than node count alone. No theorem in this run forces an arbitrary SAT algorithm to instantiate a bounded-repetition branching program, so the arbitrary-algorithm bridge remains OPEN.

\paragraph{Next target.}
Attempt symmetry-reduced SAT/ILP synthesis for a $\le61$-node read-twice BP, and extend the STORE/RECOMP experiment to the five-variable clause-presence family where each input variable occurs four times.

\begin{thebibliography}{9}
\bibitem{JR98} S.~Jukna and A.~Razborov, \emph{Neither Reading Few Bits Twice nor Reading Illegally Helps Much}, Discrete Applied Mathematics 85(3):223--238, 1998.
\bibitem{SZ00} P.~Savicky and S.~Zak, \emph{A Read-Once Lower Bound and a $(1,+k)$-Hierarchy for Branching Programs}, Theoretical Computer Science 238:347--362, 2000.
\bibitem{Woelfel02} P.~Woelfel, \emph{A Lower Bound Technique for Restricted Branching Programs and Applications}, STACS 2002, LNCS 2285:431--442.
\bibitem{Razgon25} I.~Razgon, \emph{The Splitting Power of Branching Programs of Bounded Repetition and CNFs of Bounded Width}, Discrete Applied Mathematics 360:366--381, 2025.
\bibitem{LiMcKenzie24} Y.~Li and P.~McKenzie, \emph{Perspective on Complexity Measures Targeting Read-Once Branching Programs}, Information and Computation 301:105230, 2024.
\end{thebibliography}

% END_OF_NOTE::READ-RECOMPUTE-Q4-8